\section{Resultados}

\subsection{Baseline técnico}

El holdout temporal contiene 40.687 reservas de enero a agosto de 2017, con
una tasa de cancelación de 38,70\,\%. El conjunto de desarrollo contiene
78.703 reservas de 2015 y 2016, con una tasa de 36,19\,\%.

El \texttt{DummyClassifier} asigna a todas las observaciones la probabilidad
previa de cancelación del conjunto de desarrollo. Como esta probabilidad es
inferior a 0,5, con dicho umbral predice siempre la clase mayoritaria. Su
accuracy de 0,6130 no implica capacidad predictiva: no detecta ninguna
cancelación y su ROC AUC es 0,5.

La búsqueda de regresión logística seleccionó $C=10$ y pesos de clase
balanceados. El ROC AUC medio de validación cruzada agrupada fue 0,8925, con
desviación estándar 0,0036. Sobre el holdout temporal descendió a 0,8422. La
baja dispersión entre folds indica estabilidad interna, pero la diferencia
frente a 2017 muestra que la generalización temporal es más exigente.

\begin{table}[htbp]
    \centering
    \caption{Resultados del baseline sobre el holdout temporal.}
    \label{tab:baseline-results}
    \begin{tabular}{lrrrrr}
        \toprule
        Modelo & ROC AUC & Accuracy & Precision & Recall & F1 \\
        \midrule
        Dummy prior & 0,5000 & 0,6130 & 0,0000 & 0,0000 & 0,0000 \\
        Regresión logística & 0,8422 & 0,7061 & 0,5798 & 0,8740 & 0,6971 \\
        \bottomrule
    \end{tabular}
\end{table}

Con umbral 0,5, la regresión logística detectó 13.761 de las 15.745
cancelaciones. La matriz contiene 14.967 verdaderos negativos, 9.975 falsos
positivos, 1.984 falsos negativos y 13.761 verdaderos positivos. El recall
elevado es coherente con los pesos balanceados, pero se obtiene a costa de una
cantidad considerable de falsos positivos.

\subsection{Primer lote de ingeniería de variables}

La comparación se realizó exclusivamente mediante validación cruzada agrupada
sobre 2015--2016. Eliminar \texttt{ArrivalDateWeekNumber} produjo una caída de
0,000168 respecto al baseline histórico, considerada marginal frente a la
reducción de redundancia temporal.

\begin{table}[htbp]
    \centering
    \caption{Comparación combinada del primer lote para regresión logística.}
    \label{tab:first-feature-batch}
    \begin{tabular}{lrrr}
        \toprule
        Conjunto & Features & ROC AUC CV & Desviación \\
        \midrule
        Base depurada & 24 & 0,892322 & 0,003538 \\
        Baseline histórico & 25 & 0,892490 & 0,003574 \\
        Primer lote parsimonioso & 25 & 0,893810 & 0,003613 \\
        Candidato redundante & 29 & 0,894158 & 0,003896 \\
        \bottomrule
    \end{tabular}
\end{table}

El conjunto parsimonioso superó ambas referencias en los cinco folds. El
candidato redundante lo superó solamente en 0,000348 y requirió cuatro
variables adicionales con asociaciones casi determinísticas. Se seleccionó
por tanto el conjunto parsimonioso, que conserva aproximadamente el 79\,\% de
la mejora adicional del candidato redundante sobre la base depurada y ofrece
coeficientes más interpretables. El holdout temporal no se consultó durante
esta decisión.

\subsection{Segundo lote de ingeniería de variables}

La transformación de \texttt{LeadTime} produjo la mejora principal. Agregar
\texttt{log1p(LeadTime)} al conjunto parsimonioso elevó el ROC AUC medio a
0,904913. \texttt{HasSpecialRequests} también mejoró de forma consistente y,
al combinarlo con ambas representaciones de \texttt{LeadTime}, se alcanzó
0,908554. La alternativa que reemplaza el valor continuo por bandas operativas
y agrega \texttt{HasSpecialRequests} obtuvo 0,908447, con desviación 0,004005.

\begin{table}[htbp]
    \centering
    \caption{Comparación combinada del segundo lote para regresión logística.}
    \label{tab:second-feature-batch}
    \begin{tabular}{lrrr}
        \toprule
        Configuración & Features & ROC AUC CV & Desviación \\
        \midrule
        Primer lote parsimonioso & 25 & 0,893810 & 0,003613 \\
        \texttt{LeadTime} + logaritmo & 26 & 0,904913 & 0,004321 \\
        \texttt{LeadTime} + logaritmo + solicitudes & 27 & 0,908554 & 0,004228 \\
        Bandas de anticipación + solicitudes & 26 & 0,908447 & 0,004005 \\
        \bottomrule
    \end{tabular}
\end{table}

La diferencia entre las dos mejores configuraciones fue de sólo 0,000107. Se
seleccionó la alternativa basada en bandas porque ofrece una interpretación
operativa directa, evita conservar simultáneamente dos representaciones muy
asociadas de \texttt{LeadTime} y presenta una dispersión ligeramente menor.
La mejora frente al primer lote fue de 0,014637 y se observó en los cinco
folds. El holdout temporal de 2017 permaneció reservado.

\subsection{Regresión logística mejorada}

La búsqueda seleccionó $C=1$, pesos de clase balanceados y regularización L1.
El ROC AUC medio de validación cruzada fue 0,908557, con desviación 0,005873 y
un gap train--validación de 0,003222. Sobre el holdout temporal alcanzó
0,855017, una mejora de 0,012817 frente al baseline. La caída de 0,053540 entre
CV y 2017 indica que persiste una diferencia de generalización temporal, pero
la mejora obtenida durante desarrollo se conserva fuera de muestra.

\begin{table}[htbp]
    \centering
    \caption{Comparación de Logistic Regression sobre el holdout temporal.}
    \label{tab:logistic-comparison}
    \begin{tabular}{lrrrrr}
        \toprule
        Modelo & ROC AUC & Accuracy & Precision & Recall & F1 \\
        \midrule
        Baseline & 0,842200 & 0,706100 & 0,579800 & 0,874000 & 0,697100 \\
        Mejorado & 0,855017 & 0,739032 & 0,620290 & 0,839568 & 0,713461 \\
        \bottomrule
    \end{tabular}
\end{table}

Con threshold 0,5 mejoraron Accuracy, Precision y F1, mientras que Recall se
redujo en 0,034432. Este intercambio no se optimizó en esta etapa porque ROC
AUC no depende del threshold y la política de decisión se comparará de manera
común entre los modelos.

La regularización L1 anuló 79 de las 229 columnas transformadas. Las bandas de
\texttt{LeadTime} exhibieron un patrón monotónico: la cancelación observada
aumentó desde 6,85\,\% para reservas del mismo día hasta 77,63\,\% para
anticipaciones superiores a 365 días. \texttt{HasSpecialRequests},
\texttt{TotalOfSpecialRequests} y \texttt{RequiredCarParkingSpaces} se
asociaron con menor riesgo, mientras que \texttt{PreviousCancellationRate} se
asoció con mayor riesgo.

Los mayores coeficientes positivos correspondieron a algunas categorías con
muy pocos registros. Por ejemplo, \texttt{DistributionChannel=Undefined}
apareció cinco veces y \texttt{ReservedRoomType=P}, seis. Estos valores se
consideran señales inestables producidas por baja frecuencia y no evidencia
general de importancia, aunque conservaran el signo entre folds.

\subsection{Árbol de decisión}

La importancia por permutación identificó como núcleo predictivo a
\texttt{Country}, \texttt{MarketSegment}, \texttt{PreviousCancellations},
\texttt{LeadTime}, \texttt{TotalOfSpecialRequests},
\texttt{ArrivalDateYear}, \texttt{RequiredCarParkingSpaces} y
\texttt{CustomerType}. Todas mostraron importancia positiva en los cinco
folds. \texttt{LeadTime} participó en 40 splits, confirmando la relevancia de
su relación no lineal con la cancelación.

La poda de 17 variables elevó el ROC AUC interno de 0,923749 a 0,924418 y
mejoró los cinco folds. La búsqueda final seleccionó profundidad 14,
\texttt{min\_samples\_leaf=75}, entropía y pesos balanceados. El árbol final
contiene 392 hojas y 783 nodos.

\begin{table}[htbp]
    \centering
    \caption{Comparación técnica del árbol y Logistic Regression.}
    \label{tab:decision-tree-comparison}
    \begin{tabular}{lrrrrrr}
        \toprule
        Modelo & ROC AUC CV & ROC AUC 2017 & Accuracy & Precision & Recall & F1 \\
        \midrule
        LR baseline & 0,892500 & 0,842200 & 0,706100 & 0,579800 & 0,874000 & 0,697100 \\
        LR mejorada & 0,908557 & 0,855017 & 0,739032 & 0,620290 & 0,839568 & 0,713461 \\
        Árbol podado & 0,927092 & 0,838694 & 0,765822 & 0,723456 & 0,639187 & 0,678716 \\
        \bottomrule
    \end{tabular}
\end{table}

Aunque el árbol obtuvo el mejor ROC AUC en validación cruzada, descendió a
0,838694 sobre 2017 y presentó un gap temporal de 0,088398. Con threshold 0,5
fue más preciso y alcanzó mayor Accuracy, pero omitió una proporción mayor de
cancelaciones y obtuvo menor Recall y F1 que la regresión mejorada. Se conserva
como herramienta para comprender y reducir variables, pero no como el mejor
modelo temporal. El resultado del holdout no se utilizó para una nueva
búsqueda.

\subsection{Random Forest}

Con un bosque fijo, la base depurada de 24 variables obtuvo ROC AUC medio
0,926422 frente a 0,925012 del conjunto podado y ganó cuatro de los cinco
folds. Esto indica que algunas variables poco útiles para un árbol individual
pueden contribuir a la diversidad del ensemble.

La búsqueda amplia seleccionó profundidad libre,
\texttt{max\_features=0.5} y \texttt{min\_samples\_leaf=5}. La grilla de
confirmación con 100 árboles reprodujo la selección y alcanzó ROC AUC CV
0,941260, sólo 0,000307 menos que la ejecución con 300 árboles. El ROC AUC de
entrenamiento fue 0,987653, con gap train--CV de 0,046393.

\begin{table}[htbp]
    \centering
    \caption{Resultado temporal de Random Forest.}
    \label{tab:random-forest-results}
    \begin{tabular}{lrrrrrr}
        \toprule
        Modelo & ROC AUC CV & ROC AUC 2017 & Accuracy & Precision & Recall & F1 \\
        \midrule
        LR mejorada & 0,908557 & 0,855017 & 0,739032 & 0,620290 & 0,839568 & 0,713461 \\
        Árbol podado & 0,927092 & 0,838694 & 0,765822 & 0,723456 & 0,639187 & 0,678716 \\
        Random Forest & 0,941260 & 0,842641 & 0,769312 & 0,748885 & 0,607621 & 0,670898 \\
        \bottomrule
    \end{tabular}
\end{table}

Random Forest presentó el mejor Accuracy y Precision de los modelos evaluados,
pero detectó sólo el 60,76\,\% de las cancelaciones. El ROC AUC cayó a 0,842641
en 2017, generando un gap temporal de 0,098618, y permaneció por debajo de
Logistic Regression mejorada tanto en ROC AUC temporal como en F1. La
configuración se cerró sin reajustarla a partir del holdout.

\subsection{Cancelaciones críticas}

Dentro del holdout se identificaron 348 \texttt{No-Show} y 1.250
cancelaciones ocurridas como máximo siete días antes de la llegada: 1.598
eventos críticos en total. El baseline detectó 1.179 y omitió 419, alcanzando
un recall crítico de 73,78\,\%.

El valor bruto estimado de los eventos críticos fue 540.098,62 unidades
monetarias. El modelo identificó 419.003,64 (77,58\,\%) y dejó sin identificar
121.094,98. Estos valores representan exposición bruta bajo los supuestos
definidos, no ingresos efectivamente perdidos o recuperados.

\begin{table}[htbp]
    \centering
    \caption{Evaluación operativa del baseline con umbral 0,5.}
    \label{tab:baseline-operational}
    \begin{tabular}{lr}
        \toprule
        Indicador & Resultado \\
        \midrule
        Eventos críticos & 1.598 \\
        Eventos críticos detectados & 1.179 \\
        Recall crítico & 73,78\,\% \\
        Intervenciones & 23.736 \\
        Intervenciones sobre casos no críticos & 22.557 \\
        Precisión de intervención crítica & 4,97\,\% \\
        Valor crítico capturado & 77,58\,\% \\
        \bottomrule
    \end{tabular}
\end{table}

Intervenir sobre 23.736 reservas equivale al 58,34\,\% del holdout. Una
selección aleatoria de igual tamaño detectaría en promedio aproximadamente 932
eventos críticos, frente a los 1.179 del modelo. El lift observado es 1,2647:
la regresión logística identifica alrededor de 26,5\,\% más eventos críticos
que una política aleatoria con la misma capacidad.

\subsection{Escenario económico ilustrativo}

Como comprobación de la metodología se evaluó un escenario que supone pérdida
del 100\,\% del valor bruto de cada evento crítico, costo de 10 unidades por
intervención, efectividad de 25\,\% y costo adicional nulo para los casos no
críticos. Estos parámetros son ilustrativos y no fueron observados en los
datos.

Bajo dichos supuestos, no intervenir representa 540.098,62 unidades de costo
estimado y aplicar la política del baseline representa 672.707,71. El ahorro esperado es negativo en 132.609,09 unidades
\SI{-24,55}{\percent}. Con ese costo de intervención, la efectividad mínima para
alcanzar el punto de equilibrio sería 56,65\,\%.

Por lo tanto, la política basada en el baseline y el umbral 0,5 no resulta
conveniente bajo el escenario analizado. Esta conclusión no se generaliza a
todo uso de modelos predictivos: el baseline contiene señal y supera la
selección aleatoria, pero distribuye las intervenciones sobre demasiadas
reservas. Las siguientes etapas evaluarán si nuevas variables, otros modelos y
una política de capacidad limitada permiten concentrar mejor el riesgo.
